\documentclass[12pt]{article}
\usepackage[utf8]{inputenc}
\usepackage[T1]{fontenc}
\usepackage{hyperref}
\title{Tobilung Language Resources for K-12 Curriculum Development}
\author{Compiled by Claude Sonnet 5 (claude-sonnet-5) for Andr\'as Kornai}
\date{}
\begin{document}
\maketitle

\section{Introduction}
Tobilung (ISO 639-3 tgb; Glottocode tobi1239) is an Austronesian language of Sabah, Malaysia, in the Dusunic group. It is spoken in Kota Belud, Kota Marudu and Kudat districts and in some towns. Webonary reports about 4,000 speakers, with about half of children learning it. These sources are community-oriented, so the figures are rough.

\section{Primary Educational Resources}
No school primers, textbooks or school programmes were found in the search results.

\section{Linguistic and Reference Materials}
\begin{itemize}
\item Webonary Tobilung dictionary site and introduction: \url{https://www.webonary.work/tobilung/overview/introduction/}
\item A Webonary entry page: \url{https://www.webonary.work/tobilung/en/?p=9444}
\item Glottolog entry: \url{https://glottolog.org/resource/languoid/id/tobi1239.bigmap.html}
\item Wikipedia, Dusunic languages (classification): \url{https://en.wikipedia.org/wiki/Dusunic_languages}
\item UMS paper on Tobilung traditional beliefs (komburongo), based on fieldwork in six villages: \url{https://jurcon.ums.edu.my/ojums/index.php/GA/article/download/3636/2393/13408}
\item UMS open educational resources record (relevance to Tobilung unverified): \url{https://oer.ums.edu.my/handle/oer_source_files/2509}
\end{itemize}

\section{Religious and Community Resources}
\begin{itemize}
\item Joshua Project, Tobilung in Malaysia: \url{https://joshuaproject.net/people_groups/15523/MY}
\end{itemize}
The search summary mentions Walai Tobilung (Tobilung Cultural House) in Kampung Minansad, Kota Marudu, but no page for it was returned.

\section{Digital and Audio Resources}
Webonary states that videos and literature in Tobilung exist, but no specific items or URLs were found. The Webonary dictionary is the one confirmed digital resource.

\section{Curriculum Development Context}
The language is not known to be taught in schools. No grammar or Bible translation was found. A curriculum could start from the Webonary dictionary and community partners such as the Tobilung Cultural House. Note that Lotud is a different language and its resources must not be substituted.
\end{document}